
\begin{exercise}
From the text, Exercise 3.6.
\end{exercise}

\begin{exercise}
From the text, Exercise 3.8.
\end{exercise}

\begin{exercise}
From the text, Exercise 3.9.
\end{exercise}

\begin{exercise}
Let $G$ be a connected 3-regular planar graph on 8 vertices.  How many regions does $G$ have?  Give an example of a 3-regular connected graph on 8 vertices that is not planar or explain why it must be planar.
\end{exercise}

\begin{exercise}
For the following degree sequences (list of the degrees of the vertices in a graph) determine if there exists a bipartite simple graph with these degrees.  If yes construct it, if not explain why it does not exist.\\
a. (6, 6, 4, 4, 4, 4, 4, 4, 2, 2, 2)\\ b. (4, 4, 4, 4, 4, 3, 3, 2, 2)\\ c. (5, 5, 5, 5, 5, 5, 4, 4, 4)
\end{exercise}


\begin{exercise}
What is the smallest number of colors that can be used to color the vertices of a cube so that no two adjacent vertices are colored identically?
\end{exercise}

\begin{exercise}
 Euler's formula $v-e+f=2$ holds for all connected simple planar graphs. What if a graph is not connected? Suppose a planar graph has two components. What is the value of v-e+f
 now? What if it has $k$ components?
 \end{exercise}

\begin{proofp}
Prove the chromatic number of any tree is two. Recall, a tree is a connected simple  graph with no cycles.
\end{proofp}

\begin{proofp}
Prove that if each component of a simple graph is bipartite, then the entire graph is bipartite.`
\end{proofp}

\begin{proofp}
Let $G$ be a connected simple graph on $n$ vertices.  Prove that if $n \ge 11$ then either $G$ or the complement of $G$ is not planar.
\end{proofp}

\begin{proofp}
Prove that if a simple connected planar graph has $e$ edges and $v$ vertices with $v \ge 3$ and no cycles of length 3 then $e \le 2v-4$.
\end{proofp}

\begin{proofp}
Prove Euler's formula using induction on the number of  {\bf edges} in the graph.
\end{proofp}

\begin{proofp}
Prove $G$ is regular if and only if it's complement is regular.
\end{proofp}

\begin{proofp}$\bigstar$
Prove that any simple connected planar graph must have a vertex of degree 5 or less.
\end{proofp}

\begin{proofp}$\bigstar$
Prove that a simple graph with 17 vertices and 73 edges cannot be bipartite. (Hint: Break the graph up into two vertex sets and count the number of edges within each vertex set.)
\end{proofp}

\begin{proofp} $\bigstar$
Prove Dirac's Theorem: If $G$ is a simple graph with $p$ vertices, each vertex having degree $\ge \frac{1}{2}p$, then $G$ is hamiltonian. Hint: Suppose $G$ is not hamiltonian but adding any one edge will make $G$ hamiltonian, so $G$ has a path $v_1 \rightarrow v_2 \rightarrow \cdots \rightarrow v_p$ containing all the vertices of $G$.  Show there must be a vertex adjacent to both $v_1$ and $v_p$, and use this to create a hamiltonian cycle. 
\end{proofp}

\begin{proofp}$\bigstar$
Let $G$ be a graph (not necessarily connected as in a problem above) on $n$ vertices.  Prove that if $n \ge 11$ then either $G$ or the complement of $G$ is not planar.
\end{proofp}

%\begin{exercise}
%Create a list of non-isomorphic trees on 7 vertices.  How many of the these graphs satisfy the condition: the degree of each vertex is at most 2?
%5\end{exercise}


%\begin{exercise}
%From the text, Exercise 3.6.
%\end{exercise}

%\begin{exercise}
%From the text, Exercise 3.8.
%\end{exercise}

%\begin{exercise}
%Find and draw all non-isomorphic bipartite simple graphs with exactly 5 vertices and at least one cycle.
%\end{exercise}


%\begin{proofp}
%From the text, Exercise 3.12 b) (complete matching is defined in 3.11).
%\end{proofp}


